\item Usted está trabajando para una compañía de neumáticos para automóviles. Su supervisor está estudiando los efectos de las moléculas que golpean la superficie interna del neumático debido a su movimiento térmico. Él le da los siguientes datos de un experimento reciente. Se midió el aire en el neumático de un automóvil estacionado para tener una presión manométrica de $P_1 = 1.65\text{ atm}$ en un día cuando la temperatura era de $T = 6.5^\circ\text{C}$. El automóvil fue conducido por un tiempo y luego se tomaron medidas de nuevo. La presión manométrica en el neumático era entonces de $P_2 = 1.95\text{ atm}$ y el volumen interior del neumático había aumentado en un 5.00\%. (a) Su supervisor le pide que determine por qué factor la velocidad eficaz de las moléculas de aire se ha incrementado desde la primera medición hasta la segunda. (b) También insinúa una propuesta que va a hacer para reemplazar el aire en los neumáticos con argón. ¿Cambiará esto el factor por el cual la velocidad promedio de las moléculas cambia en las condiciones descritas?
